% SPDX-License-Identifier: CC-BY-SA-4.0
% Author: Matthieu Perrin

\begingroup

\begin{frame}{Système réparti}

  \vfill
  \begin{block}{Définition -- Système réparti}
    Un système réparti est un ensemble d'entités séquentielles
    qui doivent communiquer et coopérer pour régler un
    problème commun.
    \begin{description}
    \item[Vocabulaire] distributed (anglais) = réparti (français)
    \end{description}
  \end{block}

  \vfill
  \begin{exampleblock}{Exemples}
    \begin{itemize}
    \item Entités séquentielles
      \begin{itemize}
      \item Automate définissant une séquence de \alert{pas}
      \item \structure{Threads}, \structure{processus}, ordinateur, capteur, êtres vivants...
      \end{itemize}
    \item Communication
      \begin{itemize}
      \item \structure{Mémoire partagée}, réseau, publications, gestes, phéromones...
      \end{itemize}
    \end{itemize}
  \end{exampleblock}
  \vfill  

  \pause
  \begin{alertblock}{Hypothèse simplificatrice}
    \begin{itemize}
    \item Un ensemble de \alert{$n$ processus} (ou threads) $p_1, ..., p_n$
      \begin{itemize}
      \item nombre $n$ fixe au cours d'une exécution, et connu
      \item chaque processus $p_i$ connaît son identifiant $i$
      \end{itemize}
    \end{itemize}
  \end{alertblock}
  \vfill  
\end{frame}

\endgroup
\endinput
